\pagestyle{empty}
\tight
\begin{tblr}{width=\textwidth,colspec={|Q[c,m,1.9cm]|X[l,m]|},hlines,vlines,hspan=minimal,row{1}={bg=headblue},rowsep=1pt,colsep=3pt}
\SetCell{c,m}\hfil\logo\hfil & \textbf{POLITEKNIK NEGERI MALANG}\par\textbf{JURUSAN TEKNOLOGI INFORMASI}\par\textbf{PROGRAM STUDI: D4 TEKNIK INFORMATIKA} \\
\SetCell[c=2]{c} \textbf{RENCANA PEMBELAJARAN SEMESTER (RPS)} & \\
\end{tblr}
\begin{longtblr}{width=\textwidth,colspec={|Q[l,m,1.9cm]|X[0.85,l,m]|X[1.45,l,m]|X[0.75,l,m]|X[1.15,l,m]|},hlines,vlines,hspan=minimal,rowsep=1pt,colsep=2pt,row{1,3}={bg=headgray,font=\bfseries}}
MATA KULIAH & KODE & BOBOT (SKS)/jam & SEMESTER & TGL. PENYUSUNAN \\
Metodologi Penelitian & RTI255002 & 2 SKS/4 jam & 5 & 1 Juli 2025 \\
OTORISASI & Dosen Pengembang RPS & Koordinator RMK & \SetCell[c=2]{l} Ka PRODI &  \\
 & Triana Fatmawati, S.T., M.T. & Triana Fatmawati, S.T., M.T. & \SetCell[c=2]{l}{\vspace{8mm}\par Dr. Ely Setyo Astuti, S.T., M.T.} &  \\
\SetCell[r=13]{l} Capaian Pembelajaran (CP) & \SetCell[c=4]{l,bg=headgray,font=\bfseries} Capaian Pembelajaran Lulusan Program Studi (CPL-Prodi) &  &  &  \\
 & CPL04 & \SetCell[c=3]{l} Mampu mendeskripsikan secara saintifik hasil kajian, pengembangan, dan implementasi TIK dalam bentuk skripsi, laporan, atau artikel ilmiah. &  &  \\
 & CPL10 & \SetCell[c=3]{l} Memiliki kompetensi untuk menganalisis persoalan kompleks untuk mengidentifikasi solusi bidang informatika/ilmu komputer dengan mempertimbangkan wawasan ilmu multidisiplin. &  &  \\
 & \SetCell[c=4]{l,bg=headgray,font=\bfseries} Capaian Pembelajaran mata kuliah (CPL-MK) &  &  &  \\
 & CPMK0403 & \SetCell[c=3]{l} Mahasiswa mampu merumuskan masalah penelitian, metodologi, dan menyusun proposal penelitian TI secara ilmiah. &  &  \\
 & CPMK0401 & \SetCell[c=3]{l} Mahasiswa mampu menganalisis masalah kompleks dan menghasilkan solusi rasional berbasis data yang dideskripsikan secara saintifik. &  &  \\
 & CPMK1009 & \SetCell[c=3]{l} Mampu mengidentifikasi tahapan penyelesaian ilmiah yang tepat untuk persoalan kompleks di bidang informatika. &  &  \\
 & \SetCell[c=4]{l,bg=headgray,font=\bfseries} Kemampuan akhir tiap tahapan belajar (Sub-CPMK) &  &  &  \\
 & SCPMK0403-03801 & \SetCell[c=3]{l} Mahasiswa mampu menjelaskan jenis-jenis penelitian TI, etika penelitian, dan melakukan studi literatur yang sistematis. &  &  \\
 & SCPMK0403-03802 & \SetCell[c=3]{l} Mahasiswa mampu merumuskan masalah, tujuan, manfaat, dan kerangka konseptual penelitian TI. &  &  \\
 & SCPMK0403-03803 & \SetCell[c=3]{l} Mahasiswa mampu menyusun metode penelitian (pengumpulan data, instrumen, validitas) dan penulisan proposal skripsi/LA. &  &  \\
 & SCPMK0401-03804 & \SetCell[c=3]{l} Mahasiswa mampu menganalisis data penelitian secara kuantitatif/kualitatif dan mempresentasikan hasil secara saintifik. &  &  \\
 & SCPMK1009-03805 & \SetCell[c=3]{l} Mahasiswa mampu mengidentifikasi tahapan penyelesaian ilmiah yang tepat untuk menyelesaikan persoalan kompleks di bidang informatika. &  &  \\
\SetCell[r=2]{l} Penilaian & \SetCell[c=4]{l} *Nilai CPMK didapat dari agregasi nilai SUB-CPMK yang dibebankan &  &  &  \\
 & \SetCell[c=4]{l}{\tiny\begin{tblr}{width=\linewidth,colspec={|Q[c,m,0.1\linewidth]|Q[c,m,0.16\linewidth]|X[l,m]|X[l,m]|X[l,m]|X[l,m]|X[l,m]|X[l,m]|Q[c,m,0.08\linewidth]|},hlines,vlines,hspan=minimal,rowsep=1pt,colsep=0.7pt,row{1,2}={bg=headgray,font=\bfseries}}
Kode CPMK & Kode SCPMK & \SetCell[c=6]{c} Bobot per Bentuk Penilaian & & & & & & Total Bobot \\
 & & Tugas & Kuis 1 & UTS & Kuis 2 & PBL & UAS & \\
\SetCell[r=3]{c} CPMK0403 & SCPMK0403-03801 & 4.50\%: studi literatur, jenis penelitian, etika & 5\%: konsep metodologi penelitian & & & 5\%: kontribusi proposal & & 14.50\% \\
 & SCPMK0403-03802 & 3\%: rumusan masalah, kerangka konseptual & & & & 15\%: proposal bab 1 dan 2 & & 18\% \\
 & SCPMK0403-03803 & 5.50\%: metode penelitian, instrumen, penulisan & & & & 25\%: proposal bab 3 & 8\%: ujian metodologi & 38.50\% \\
\SetCell[r=1]{c} CPMK0401 & SCPMK0401-03804 & 2\%: analisis data, presentasi & & & 5\%: analisis data dan pemodelan & 10\%: presentasi proposal & 9\%: ujian akhir & 26\% \\
\SetCell[r=1]{c} CPMK1009 & SCPMK1009-03805 & & & & & & 3\%: identifikasi tahapan penyelesaian ilmiah & 3\% \\
\SetCell[c=8]{r}\textbf{Total Per Penilaian} & & & & & & & & \textbf{100\%} \\
\end{tblr}} &  &  &  \\
Diskripsi Singkat Mata Kuliah & \SetCell[c=4]{l} Mata kuliah Metodologi Penelitian membangun kemampuan mahasiswa dalam merencanakan, merancang, dan melaksanakan penelitian TI secara ilmiah, mencakup studi literatur, identifikasi masalah, penyusunan metode, analisis data, dan penulisan proposal penelitian/skripsi. &  &  &  \\
Bahan Kajian / Materi Pembelajaran & \SetCell[c=4]{l}\begin{enumerate}\item Jenis penelitian dan etika ilmiah\item Studi literatur dan identifikasi masalah\item Metode penelitian dan pengumpulan data\item Analisis data dan pemodelan sistem\item Penulisan proposal penelitian\end{enumerate} &  &  &  \\
\SetCell[r=2]{l} Pustaka & \SetCell[c=4]{l} \textbf{Utama:}\begin{enumerate}\item Sugiyono. 2019. Metode Penelitian Kuantitatif, Kualitatif, dan R\&D. Alfabeta.\item Creswell, J.W. 2014. Research Design: Qualitative, Quantitative, and Mixed Methods. SAGE.\item Panduan Penulisan Laporan Akhir/Skripsi JTI Polinema.\end{enumerate} &  &  &  \\
 & \SetCell[c=4]{l} \textbf{Pendukung:} Materi LMS, artikel ilmiah dari jurnal nasional/internasional, dan jobsheet praktikum. &  &  &  \\
Media Pembelajaran & \SetCell[c=2]{l} \textbf{Software:} Word processor, spreadsheet, reference manager (Mendeley/Zotero), LMS. &  & \SetCell[c=2]{l} \textbf{Hardware:} Komputer/laptop, proyektor. &  \\
Nama Dosen Pengampu & \SetCell[c=4]{l} Tim Pengajar Program Studi D4 TI &  &  &  \\
Matakuliah Syarat & \SetCell[c=4]{l} RTI254008 Statistik Komputasi &  &  &  \\
\end{longtblr}
\begin{longtblr}{width=\textwidth,colspec={|Q[c,m,0.06\textwidth]|X[1.35,l,m]|X[1.08,l,m]|X[1.16,l,m]|X[0.7,l,m]|X[1.24,l,m]|X[1.06,l,m]|X[0.98,l,m]|Q[c,m,0.05\textwidth]|},hlines,vlines,hspan=minimal,rowhead=2,row{1,2}={bg=headgreen,font=\bfseries},rowsep=1pt,colsep=1.5pt}
Mgg.\par Ke & Kemampuan Akhir Yang Direncanakan (Sub-CP-MK) & Bahan kajian (Materi Pembelajaran) & Bentuk dan Metode Pembelajaran & Estimasi Waktu & Pengalaman Belajar Mahasiswa & Kriteria \& Bentuk Penilaian & Indikator Penilaian & Bobot Penilaian (\%) \\
(1) & (2) & (3) & (4) & (5) & (6) & (7) & (8) & (9) \\
1 & SCPMK0403-03801 & Pengantar Metodologi Penelitian TI. & Modalitas: Blended Learning. Bentuk: Luring/Daring. Metode: Ceramah, diskusi, penulisan ilmiah, dan case study. & 1 x 2 x 50' tatap muka; 1 x 2 x 50' tugas/praktik mandiri. & Mahasiswa mengerjakan aktivitas terarah untuk pengantar metodologi penelitian TI dan menunjukkan evidence hasil belajar. & Bentuk penilaian: Pengantar Metodologi Penelitian TI. Mengacu rubrik RTM. & Ketepatan konsep; kualitas artefak/praktik; validasi dan dokumentasi. & 1.5 \\
2 & SCPMK0403-03801 & Jenis-jenis Penelitian dan Etika Penelitian. & Modalitas: Blended Learning. Bentuk: Luring/Daring. Metode: Ceramah, diskusi, penulisan ilmiah, dan case study. & 1 x 2 x 50' tatap muka; 1 x 2 x 50' tugas/praktik mandiri. & Mahasiswa mengerjakan aktivitas terarah untuk jenis-jenis penelitian dan etika penelitian dan menunjukkan evidence hasil belajar. & Bentuk penilaian: Jenis-jenis Penelitian dan Etika Penelitian. Mengacu rubrik RTM. & Ketepatan konsep; kualitas artefak/praktik; validasi dan dokumentasi. & 1.5 \\
3 & SCPMK0403-03801 & Studi Literatur dan Review Artikel Ilmiah. & Modalitas: Blended Learning. Bentuk: Luring/Daring. Metode: Ceramah, diskusi, penulisan ilmiah, dan case study. & 1 x 2 x 50' tatap muka; 1 x 2 x 50' tugas/praktik mandiri. & Mahasiswa mengerjakan aktivitas terarah untuk studi literatur dan review artikel ilmiah dan menunjukkan evidence hasil belajar. & Bentuk penilaian: Studi Literatur dan Review Artikel Ilmiah. Mengacu rubrik RTM. & Ketepatan konsep; kualitas artefak/praktik; validasi dan dokumentasi. & 1.5 \\
4 & SCPMK0403-03802 & Identifikasi dan Rumusan Masalah. & Modalitas: Blended Learning. Bentuk: Luring/Daring. Metode: Ceramah, diskusi, penulisan ilmiah, dan case study. & 1 x 2 x 50' tatap muka; 1 x 2 x 50' tugas/praktik mandiri. & Mahasiswa mengerjakan aktivitas terarah untuk identifikasi dan rumusan masalah dan menunjukkan evidence hasil belajar. & Bentuk penilaian: Identifikasi dan Rumusan Masalah. Mengacu rubrik RTM. & Ketepatan konsep; kualitas artefak/praktik; validasi dan dokumentasi. & 1.5 \\
5 & SCPMK0403-03801 & Kuis 1 / Quiz 1. & Modalitas: Blended Learning. Bentuk: Luring/Daring. Metode: Ceramah, diskusi, penulisan ilmiah, dan case study. & 1 x 2 x 50' tatap muka; 1 x 2 x 50' tugas/praktik mandiri. & Mahasiswa mengerjakan kuis 1 untuk mengukur pemahaman konsep metodologi penelitian. & Bentuk penilaian: Kuis 1. Mengacu rubrik RTM. & Ketepatan konsep; kualitas artefak/praktik; validasi dan dokumentasi. & 5 \\
6 & SCPMK0403-03802 & Tujuan, Manfaat, dan Hipotesis Penelitian. & Modalitas: Blended Learning. Bentuk: Luring/Daring. Metode: Ceramah, diskusi, penulisan ilmiah, dan case study. & 1 x 2 x 50' tatap muka; 1 x 2 x 50' tugas/praktik mandiri. & Mahasiswa mengerjakan aktivitas terarah untuk tujuan, manfaat, dan hipotesis penelitian dan menunjukkan evidence hasil belajar. & Bentuk penilaian: Tujuan, Manfaat, dan Hipotesis Penelitian. Mengacu rubrik RTM. & Ketepatan konsep; kualitas artefak/praktik; validasi dan dokumentasi. & 1.5 \\
7 & SCPMK0403-03803 & Metode Pengumpulan Data (Kuantitatif dan Kualitatif). & Modalitas: Blended Learning. Bentuk: Luring/Daring. Metode: Ceramah, diskusi, penulisan ilmiah, dan case study. & 1 x 2 x 50' tatap muka; 1 x 2 x 50' tugas/praktik mandiri. & Mahasiswa mengerjakan aktivitas terarah untuk metode pengumpulan data dan menunjukkan evidence hasil belajar. & Bentuk penilaian: Metode Pengumpulan Data. Mengacu rubrik RTM. & Ketepatan konsep; kualitas artefak/praktik; validasi dan dokumentasi. & 1.5 \\
8 & SCPMK0403-03801--SCPMK0403-03802 & UTS (Proposal Penelitian Awal). & Modalitas: Blended Learning. Bentuk: Luring/Daring. Metode: Ceramah, diskusi, penulisan ilmiah, dan case study. & 1 x 2 x 50' tatap muka; 1 x 2 x 50' tugas/praktik mandiri. & Mahasiswa menyerahkan dan mempresentasikan draft proposal penelitian awal (bab 1 dan 2). & Bentuk penilaian: UTS Proposal Penelitian Awal. Mengacu rubrik RTM. & Kelengkapan proposal; kualitas argumen; konsistensi masalah dan metode. & 20 \\
9 & SCPMK0403-03803 & Instrumen Penelitian, Validitas, dan Reliabilitas. & Modalitas: Blended Learning. Bentuk: Luring/Daring. Metode: Ceramah, diskusi, penulisan ilmiah, dan case study. & 1 x 2 x 50' tatap muka; 1 x 2 x 50' tugas/praktik mandiri. & Mahasiswa mengerjakan aktivitas terarah untuk instrumen penelitian, validitas, dan reliabilitas dan menunjukkan evidence hasil belajar. & Bentuk penilaian: Instrumen Penelitian. Mengacu rubrik RTM. & Ketepatan konsep; kualitas artefak/praktik; validasi dan dokumentasi. & 1.5 \\
10 & SCPMK0401-03804 & Pengolahan Data dan Analisis Statistik. & Modalitas: Blended Learning. Bentuk: Luring/Daring. Metode: Ceramah, diskusi, penulisan ilmiah, dan case study. & 1 x 2 x 50' tatap muka; 1 x 2 x 50' tugas/praktik mandiri. & Mahasiswa mengerjakan aktivitas terarah untuk pengolahan data dan analisis statistik dan menunjukkan evidence hasil belajar. & Bentuk penilaian: Pengolahan Data dan Analisis Statistik. Mengacu rubrik RTM. & Ketepatan konsep; kualitas artefak/praktik; validasi dan dokumentasi. & 1.5 \\
11 & SCPMK0403-03803 & Pemodelan Sistem (SDLC, Agile, Waterfall). & Modalitas: Blended Learning. Bentuk: Luring/Daring. Metode: Ceramah, diskusi, penulisan ilmiah, dan case study. & 1 x 2 x 50' tatap muka; 1 x 2 x 50' tugas/praktik mandiri. & Mahasiswa mengerjakan aktivitas terarah untuk pemodelan sistem dan menunjukkan evidence hasil belajar. & Bentuk penilaian: Pemodelan Sistem. Mengacu rubrik RTM. & Ketepatan konsep; kualitas artefak/praktik; validasi dan dokumentasi. & 1.5 \\
12 & SCPMK0401-03804 & Kuis 2 / Quiz 2. & Modalitas: Blended Learning. Bentuk: Luring/Daring. Metode: Ceramah, diskusi, penulisan ilmiah, dan case study. & 1 x 2 x 50' tatap muka; 1 x 2 x 50' tugas/praktik mandiri. & Mahasiswa mengerjakan kuis 2 untuk mengukur pemahaman analisis data dan pemodelan sistem. & Bentuk penilaian: Kuis 2. Mengacu rubrik RTM. & Ketepatan konsep; kualitas artefak/praktik; validasi dan dokumentasi. & 5 \\
13 & SCPMK0403-03803 & Teknik Penulisan Karya Ilmiah (Struktur, Sitasi, Referensi). & Modalitas: Blended Learning. Bentuk: Luring/Daring. Metode: Ceramah, diskusi, penulisan ilmiah, dan case study. & 1 x 2 x 50' tatap muka; 1 x 2 x 50' tugas/praktik mandiri. & Mahasiswa mengerjakan aktivitas terarah untuk teknik penulisan karya ilmiah dan menunjukkan evidence hasil belajar. & Bentuk penilaian: Teknik Penulisan Karya Ilmiah. Mengacu rubrik RTM. & Ketepatan konsep; kualitas artefak/praktik; validasi dan dokumentasi. & 0.5 \\
14 & SCPMK0403-03803 & Penulisan Proposal Skripsi/LA: Bab 1 dan 2. & Modalitas: Blended Learning. Bentuk: Luring/Daring. Metode: Ceramah, diskusi, penulisan ilmiah, dan case study. & 1 x 2 x 50' tatap muka; 1 x 2 x 50' tugas/praktik mandiri. & Mahasiswa mengerjakan aktivitas terarah untuk penulisan proposal bab 1 dan 2 dan menunjukkan evidence hasil belajar. & Bentuk penilaian: Penulisan Proposal Bab 1 dan 2. Mengacu rubrik RTM. & Ketepatan konsep; kualitas artefak/praktik; validasi dan dokumentasi. & 0.5 \\
15 & SCPMK0403-03803 & Penulisan Proposal Skripsi/LA: Bab 3. & Modalitas: Blended Learning. Bentuk: Luring/Daring. Metode: Ceramah, diskusi, penulisan ilmiah, dan case study. & 1 x 2 x 50' tatap muka; 1 x 2 x 50' tugas/praktik mandiri. & Mahasiswa mengerjakan aktivitas terarah untuk penulisan proposal bab 3 dan menunjukkan evidence hasil belajar. & Bentuk penilaian: Penulisan Proposal Bab 3. Mengacu rubrik RTM. & Ketepatan konsep; kualitas artefak/praktik; validasi dan dokumentasi. & 0.5 \\
16 & SCPMK0403-03801--SCPMK0401-03804 & Presentasi Proposal Penelitian. & Modalitas: Blended Learning. Bentuk: Luring/Daring. Metode: Ceramah, diskusi, penulisan ilmiah, dan case study. & 1 x 2 x 50' tatap muka; 1 x 2 x 50' tugas/praktik mandiri. & Mahasiswa mempresentasikan proposal penelitian lengkap dan mempertahankan keputusan metodologis. & Bentuk penilaian: PBL Presentasi Proposal Penelitian. Mengacu rubrik RTM. & Kualitas proposal; ketepatan metodologi; kemampuan defensi. & 35 \\
17 & SCPMK0403-03803, SCPMK0401-03804, SCPMK1009-03805 & UAS / Final Exam. & Modalitas: Blended Learning. Bentuk: Luring/Daring. Metode: Ceramah, diskusi, penulisan ilmiah, dan case study. & 1 x 2 x 50' tatap muka; 1 x 2 x 50' tugas/praktik mandiri. & Mahasiswa mengerjakan ujian akhir semester mencakup seluruh materi metodologi penelitian, termasuk identifikasi tahapan penyelesaian ilmiah untuk persoalan kompleks bidang informatika. & Bentuk penilaian: UAS Metodologi Penelitian. Mengacu rubrik RTM. & Ketepatan konsep; kemampuan analisis; penulisan ilmiah; ketepatan identifikasi tahapan penyelesaian ilmiah. & 20 \\
\end{longtblr}
